\subsection{Why small constraint entries do not bound the average polynomial}

\begin{proposition}[P-F38: an exponential obstruction for bounded-entry vector partitions]
\label{prop:boundedentryaverage}
For every integer $n\ge24$, put
\[
 M=2^n,\qquad r=2\lfloor n^3/2\rfloor.
 \tag{P-F38a}
\]
There is a rational polytope $P_n$ of dimension $r$ in the
unit box $[0,1]^{r+n+1}$, defined by nonnegativity and $n+1$
independent equations with integer coefficients of absolute
value at most two and right-hand sides zero or one, such that
its Ehrhart quasipolynomial $F_n(s)=\#(sP_n\cap\mathbb Z^{r+n+1})$
has average constituent polynomial
\[
 B_n^{\mathrm{ex}}(X)=\frac1L\sum_{a=0}^{L-1}q_a(X)
      =\sum_{j=0}^{r}b_jX^j
\]
whose coefficient norm satisfies
\[
 \sum_j|b_j|\ \ge\ |b_0|\ \ge\ 2^{n^4/2}.
 \tag{P-F38b}
\]
Here $L$ is any quasipolynomial period. The polytope, its
integral lattice, and the coefficient lower bound are explicit.
The same counting function has an equivalent unit-box
presentation with $r+2n+1$ variables and $2n+1$ independent
equations in which every matrix entry is in $\{0,1,2\}$,
every right-hand side is one, and each column sum is at most
three.
Consequently, a bound of the form
$2^{C n^3\log_2(n+2)}$ for all such average-polynomial norms
does not follow from the dimension, the number of equations,
small constraint entries, or containment in a unit box.
This is not an assertion about the full cubic weight system
or its Weyl alternating sum in P-F34.
\end{proposition}

\begin{proof}
Use coordinates
$x_1,\ldots,x_r,y_0,\ldots,y_n$. Define $P_n$ by
\[
 x_j\ge0,\quad y_i\ge0,\qquad
 y_{i+1}=2y_i\ (0\le i<n),\qquad
 \sum_{j=1}^{r}x_j+y_n=1.
 \tag{P-F38c}
\]
The $n$ chain equations are independent, and the last equation
is independent of them because it has a nonzero $x_1$
coefficient. Positive solutions exist, so the dimension is
$(r+n+1)-(n+1)=r$. Every coordinate is at most one:
$x_j\le1$ and $0\le y_i\le y_n\le1$.

For every nonnegative integer $s$, the integer solutions of
the dilated equations are in bijection with
\[
 x_1,\ldots,x_r,c\in\mathbb Z_{\ge0},
 \qquad \sum_{j=1}^{r}x_j+Mc=s.
 \tag{P-F38d}
\]
The forward map takes $c=y_0$; the inverse sets $y_i=2^ic$.
Both maps preserve the integer lattice exactly. Thus this is
not merely a rational equivalence of polytopes.

For the nonnegative-matrix presentation, introduce
$z_0,\ldots,z_{n-1}\ge0$ and replace the $n$ chain equations by
\[
 y_{i+1}+z_i=1,\qquad 2y_i+z_i=1\quad(0\le i<n),
 \qquad \sum_{j=1}^{r}x_j+y_n=1.
 \tag{P-F38d'}
\]
Subtracting the two equations in each pair recovers the
chain. Conversely the chain uniquely gives
$z_i=1-y_{i+1}\in[0,1]$; after dilation this becomes
$z_i=s-y_{i+1}$. Thus adjoining or removing the $z_i$
preserves integer points exactly for every $s$.
These $2n+1$ equations are independent. A linear relation
among their rows has zero coefficient on the last row
by inspecting an $x_j$ column. Each $z_i$ column makes
the two row coefficients in its pair opposite.
The $y_0$ column then kills the first pair, and the
$y_1,y_2,\ldots,y_{n-1}$ columns kill all successive
pairs inductively. The dimension is still $r$.
Every entry is $0$, $1$, or $2$. An internal $y_i$
column has sum three, and all other column sums are
at most two. All coordinates remain in the unit box.

Summing over all these integers gives
\[
 \sum_{s\ge0}F_n(s)t^s
          =\frac1{(1-t)^r(1-t^M)}.
 \tag{P-F38e}
\]
The rational function is proper and all its poles are roots
of unity. Its partial fractions therefore give a global
quasipolynomial, with degree $r$ and a period dividing $M$.
The Fourier argument of P-F34 identifies the average
polynomial with the part belonging to the pole at one.

We prove a quantitative coefficient estimate for (P-F38e)
for every even $r\ge4$ and $M\ge128r$. Write
\[
 T(z)=\frac{z}{1-e^{-z}},\qquad
 T(z)^r=\sum_{j\ge0}a_jz^j,\qquad a_0=1.
 \tag{P-F38f}
\]
The constant coefficient $b_0$ of the average polynomial
is the coefficient of $z^{-1}$ in
$((1-e^{-z})^r(1-e^{-Mz}))^{-1}$.
Indeed
\[
 \sum_{s\ge0}s^je^{-zs}
       =\left(-\frac{d}{dz}\right)^j\frac1{1-e^{-z}}
       =j!z^{-j-1}+\text{a function holomorphic at }z=0.
 \tag{P-F38g}
\]
For $j=0$ the same formula holds with the holomorphic
constant term $1/2$. Terms belonging to poles different
from one are holomorphic near $z=0$, so (P-F38g)
extracts exactly $b_0$.

Let $B_k$ denote the Bernoulli numbers with $B_1=-1/2$.
Their generating series gives
\[
 \frac1{1-e^{-Mz}}
 =\frac1{Mz}+\frac12+
      \sum_{\substack{k\ge2\\ k\ \mathrm{even}}}
                     \frac{B_k}{k!}M^{k-1}z^{k-1}.
\]
Consequently
\[
 b_0=\frac{a_r}{M}+\frac{a_{r-1}}2
       +\sum_{\substack{2\le k\le r\\ k\ \mathrm{even}}}
                 \frac{B_k}{k!}M^{k-1}a_{r-k}.
 \tag{P-F38h}
\]
The leading term as a Laurent polynomial in $M$ is
\[
             L_r=\frac{B_r}{r!}M^{r-1}.
 \tag{P-F38i}
\]
We now bound every other term relative to it.

For $|z|\le1/2$, the exponential series gives
\[
 |(1-e^{-z})-z|
       \le\sum_{k\ge2}\frac{|z|^k}{k!}\le |z|^2.
\]
For example the last bound follows from
$\sum_{k\ge2}u^k/k!\le u^2/(2(1-u))\le u^2$
when $0\le u\le1/2$.
On $|z|=1/r$ this yields
\[
 |T(z)|^r\le(1-1/r)^{-r}\le e^2<8.
\]
The middle inequality follows from
$-\log(1-u)\le2u$ for $0\le u\le1/2$.
Cauchy's coefficient estimate therefore proves
\[
                         |a_j|\le8r^j
                         \quad\hbox{for every }j\ge0.
 \tag{P-F38j}
\]

The even Bernoulli-number formula is
\[
 \frac{|B_k|}{k!}=\frac{2\zeta(k)}{(2\pi)^k}
           \qquad(k\ge2\ \hbox{even}).
 \tag{P-F38k}
\]
It is NIST DLMF equation~25.6.2, equivalently equation~24.8.1
at $x=0$. This exact identity also follows by integrating the
Fourier coefficients of the Bernoulli polynomials by parts.
In particular $\zeta(k)<2$ and $\zeta(r)>1$, so
\[
 \frac{|B_k|/k!}{|B_r|/r!}<2(2\pi)^{r-k}
        \quad(2\le k\le r-2,\ k\ \hbox{even}),\qquad
 |L_r|\ge\frac{2M^{r-1}}{(2\pi)^r}.
\]
The bound $\zeta(k)<2$ needs no asymptotic estimate:
$\sum_{\ell\ge2}\ell^{-k}\le\sum_{\ell\ge2}\ell^{-2}<1$.

Set $\eta=2\pi r/M$. The assumption $M\ge128r$ and
$2\pi<8$ give $\eta<1/16$. For the nonleading Bernoulli
terms in (P-F38h), put $j=r-k$. Equations (P-F38j)--(P-F38k)
give the total relative bound
\[
 \sum_{\substack{2\le k\le r-2\\ k\ \mathrm{even}}}
 \frac{|B_k|M^{k-1}|a_{r-k}|}{k!\,|L_r|}
       \le16\sum_{\substack{j\ge2\\ j\ \mathrm{even}}}\eta^j
       =\frac{16\eta^2}{1-\eta^2}
       <\frac{16}{255}.
 \tag{P-F38l}
\]
The two remaining terms satisfy
\[
 \frac{|a_{r-1}|/2}{|L_r|}
        \le4\pi\eta^{r-1}<16\eta^{r-1},
 \qquad
 \frac{|a_r|/M}{|L_r|}\le4\eta^r.
 \tag{P-F38m}
\]
Since $r\ge4$, the sum of the three relative bounds is
less than
\[
        \frac{16}{255}+\frac{16}{16^3}
                          +\frac4{16^4}<\frac12.
\]
There is thus no cancellation of even half the leading
Bernoulli term. We have proved
\[
              |b_0|\ge\frac12|L_r|
                         \ge\frac{M^{r-1}}{(2\pi)^r}.
 \tag{P-F38n}
\]

For the values (P-F38a), $M\ge128r$ for every $n\ge24$.
At $n=24$, $2^{24}>128\cdot24^3$; thereafter the ratio
$2^n/n^3$ increases, because
$2(n/(n+1))^3>1$ for $n\ge4$.
Also $r\ge n^3-1$ and $\log_2(2\pi)<3$, whence
\[
 \begin{split}
 \log_2|b_0|
 &\ge (r-1)n-r\log_2(2\pi)\\
 &>r(n-3)-n\\
 &\ge n^4-3n^3-2n+3
 >n^4/2\qquad(n\ge24).
 \end{split}
 \tag{P-F38o}
\]
This proves the stated coefficient-norm lower bound.
For each fixed $C$, $n^4/2>Cn^3\log_2(n+2)$ eventually,
so the proposed universal smaller bound is impossible
on this class of bounded-entry systems.
\end{proof}

\paragraph{The two-dimensional warning and the remaining cubic gate.}
Already the triangle
$\operatorname{conv}\{(0,0),(1,0),(0,1/M)\}$ lies inside
the unit square and has Ehrhart series
$1/((1-t)^2(1-t^M))$. Writing $s=Mk+a$, $0\le a<M$,
its exact count is
\[
 \sum_{c=0}^{k}(s-Mc+1)
   =\frac{s^2}{2M}+
       \left(\frac12+\frac1M\right)s+
       1+\frac a2-\frac aM-\frac{a^2}{2M}.
 \tag{P-F38p}
\]
Thus its average constituent constant is
$(M^2+6M+5)/(12M)$, unbounded while the containing box
is fixed. Formula (P-F38c) shows that encoding a large
denominator by small constraint entries does not repair
the issue, and (P-F38b) gives the quantitative
$2^{\Omega(n^4)}$ obstruction at cubic-scale dimension.
If the nonnegative presentation is parameterized instead by
its row rank $N=2n+1$, the stronger displayed bound in
(P-F38o) gives
\[
 \liminf_{n\to\infty}
   \frac{\log_2\bigl(\sum_j|b_j|\bigr)}{(2n+1)^4}
       \ge\frac1{16}>\frac1{24}.
 \tag{P-F38q}
\]
This follows by dividing $n^4-3n^3-2n+3$ by $(2n+1)^4$
and taking its limit. Thus the sharper constant threshold
in (P-F34u) also cannot be justified from these general
features alone.

This example is an ordinary vector-partition problem,
even in the nonnegative-matrix presentation (P-F38d').
It does not have the full list of degree-three
monomial weights and the Weyl alternating numerator
of the actual cubic invariant ring. Those specific
features would have to supply an additional cancellation
or estimate if one seeks the smaller $U_n$ bound in P-F34.
No upper or lower bound on that actual $U_n$ is inferred.
The primary identity used above is available at
\url{https://dlmf.nist.gov/25.6.E2}.
